\documentclass[conference]{IEEEtran}
\IEEEoverridecommandlockouts

\usepackage{cite}
\usepackage{amsmath,amssymb,amsfonts}
\usepackage{algorithmic}
\usepackage{algorithm}
\usepackage{graphicx}
\usepackage{textcomp}
\usepackage{xcolor}
\usepackage{booktabs}
\usepackage{multirow}
\usepackage{url}
\usepackage{subcaption}

\def\BibTeX{{\rm B\kern-.05em{\sc i\kern-.025em b}\kern-.08em
    T\kern-.1667em\lower.7ex\hbox{E}\kern-.125emX}}

\begin{document}

\title{ClearDrive AI: A Real-Time Multi-Hazard Optical Restoration and Kinematic Collision Forecasting Framework for Edge ADAS\\
\thanks{This research was conducted at the Department of CSE-AIML, Chennai Institute of Technology, Chennai, India.}
}

\author{
\IEEEauthorblockN{\textit{Shivanesh V}}
\IEEEauthorblockA{\textit{Department of CSE-AIML} \\
\textit{Chennai Institute of Technology} \\
\textit{Chennai, Tamil Nadu} \\
\textit{shivaneshv.aiml2025@citchennai.net}}
\and
\IEEEauthorblockN{\textit{Gowri R}}
\IEEEauthorblockA{\textit{Department of CSE-AIML} \\
\textit{Chennai Institute of Technology} \\
\textit{Chennai, Tamil Nadu} \\
\textit{gowri.r@citchennai.net}}
\and
\IEEEauthorblockN{\textit{Karthikeyan P}}
\IEEEauthorblockA{\textit{Department of CSE-AIML} \\
\textit{Chennai Institute of Technology} \\
\textit{Chennai, Tamil Nadu} \\
\textit{karthikeyan.p@citchennai.net}}
\and
\IEEEauthorblockN{\textit{Raj Karkee J}}
\IEEEauthorblockA{\textit{Department of CSE-AIML} \\
\textit{Chennai Institute of Technology} \\
\textit{Chennai, Tamil Nadu} \\
\textit{rajkarkeej.cse@citchennai.net}}
\and
\IEEEauthorblockN{\textit{Revathi Devi M}}
\IEEEauthorblockA{\textit{Department of CSE-AIML} \\
\textit{Chennai Institute of Technology} \\
\textit{Chennai, Tamil Nadu} \\
\textit{revathidevim.maths@citchennai.net}}
\and
\IEEEauthorblockN{\textit{Akila Gandhi GS}}
\IEEEauthorblockA{\textit{Department of CSE-AIML} \\
\textit{Chennai Institute of Technology} \\
\textit{Chennai, Tamil Nadu} \\
\textit{akilagandhigs@citchennai.net}}
}

\maketitle

\begin{abstract}
Adverse atmospheric conditions---including dense aerosol fog, torrential monsoon precipitation, zero-light rural corridors, and oncoming high-beam glare---severely compromise monocular vision sensors in Advanced Driver Assistance Systems (ADAS), inducing dangerous perceptual blindness and elevated vehicular collision rates. While modern deep learning restoration models achieve notable synthetic benchmarks, their prohibitive memory footprints and non-deterministic latencies preclude deployment on resource-constrained automotive edge microprocessors. This paper presents \textbf{ClearDrive AI}, a unified, deterministic, and edge-resilient ADAS framework executing real-time optical restoration, monocular kinematic telemetry, and in-cabin biometric driver monitoring on low-cost hardware. We formulate an $\mathcal{O}(N)$ linear-time optical transmission pipeline combining an adaptive foliage-aware Dark Channel Prior (DCP) with He's Fast Guided Filter, permanently eliminating white-clipping halo artifacts and temporal sky flickering via exponential moving average (EMA) airlight tracking. We couple an optical specular knee compression curve ($>190$ digital intensity rolloff) with Retinex dynamic range expansion to counteract blinding glare and pitch darkness. Furthermore, a monocular ground-plane pinhole geometry model projects YOLOv8n detections into 3D metric distances, deriving dynamic closing velocities, $+1.5\,\text{s}$ forward trajectory cut-in vectors, and dynamic Time-to-Collision ($TTC$) alerts. Benchmarked across 12,500 frames from RESIDE, Foggy Cityscapes, and real-world Indian highway road tests (including the 36-hairpin Ooty Ghat road), ClearDrive AI sustains \textbf{32.4 FPS} on a quad-core ARM Cortex-A72 (Raspberry Pi 4B) within an \$82 (₹6,800) Bill of Materials, achieving 22.6\,dB PSNR, 0.86 SSIM, and expanding forward hazard identification range by 340\% under severe fog.
\end{abstract}

\begin{IEEEkeywords}
Advanced Driver Assistance Systems (ADAS), Dark Channel Prior (DCP), Fast Guided Filter, Atmospheric Dehazing, Monocular Depth, Time-to-Collision (TTC), Edge Computing, Autonomous Vehicles.
\end{IEEEkeywords}

\section{Introduction}
Vehicular transportation forms the vital circulatory backbone of modern economic infrastructure. However, road traffic fatalities remain an alarming global public health and engineering crisis. According to official annual reports published by the Ministry of Road Transport and Highways (MoRTH) \cite{morth2024report}, over 460,000 traffic accidents occur annually in India alone, resulting in upwards of 168,000 fatalities and permanent disabling injuries to over 440,000 citizens. A disproportionate fraction of these catastrophic collisions occur under adverse environmental conditions: dense radiation and advection aerosol fog across northern expressways (e.g., Yamuna Expressway, NH-44) and high-altitude mountain corridors (such as the 36-hairpin Ooty Ghat Road, SH-15), blinding torrential monsoon downpours, unlit rural highways, and specular glare from aftermarket high-intensity LED headlights.

\subsection{Atmospheric Degradation \& Optical Physics}
Atmospheric scattering is an ill-posed inverse optical problem governed by the physical interaction between electromagnetic photons and suspended particulate matter. In a turbid medium, optical scattering is categorized into:
\begin{enumerate}
    \item \textbf{Rayleigh Scattering:} Occurs when the particle radius $r \ll \lambda$ (wavelength of incident light), governing fine atmospheric gaseous molecules and causing standard chromatic dispersion.
    \item \textbf{Mie Scattering:} Dominates when particle radius is comparable to or larger than optical wavelengths ($r \approx \lambda$). In automotive operating design domains (ODD), suspended water droplets, aerosol mists, and industrial smogs exhibit droplet radii between $0.5\,\mu\text{m}$ and $20\,\mu\text{m}$.
\end{enumerate}
Under Mie scattering, incident photons undergo intense forward and angular dispersion, attenuating direct transmission from scene objects while scattering ambient skylight into the camera line-of-sight. The classical physical formulation governing this phenomenon is Koschmieder's Law \cite{koschmieder1924theorie}:
\begin{equation}
\mathbf{I}(x) = \mathbf{J}(x) \cdot t(x) + \mathbf{A}(1 - t(x))
\label{eq:koschmieder}
\end{equation}
where $\mathbf{I}(x)$ is the observed sensor image at pixel coordinate $x$, $\mathbf{J}(x)$ is the true unattenuated scene radiance, $\mathbf{A}$ is the global atmospheric airlight vector, and $t(x)$ is the medium transmission map. Under homogeneous atmospheric turbidity with extinction coefficient $\beta$, transmission decays exponentially with physical scene depth $d(x)$ according to the Beer-Lambert law:
\begin{equation}
t(x) = \exp(-\beta \cdot d(x))
\label{eq:beer_lambert}
\end{equation}
Because RGB sensors record 3 color channels while the image formation model involves 4 unknowns per pixel (3 radiance channels plus scalar transmission $t(x)$), single-image optical recovery is fundamentally under-constrained and mathematically ill-posed.

\subsection{Vulnerability of Contemporary Automotive Sensors}
Modern ADAS and autonomous vehicle (AV) perception suites rely on a combination of physical sensors, each presenting critical failure modes when operating in severe tropical and developing-nation driving environments:
\begin{enumerate}
    \item \textbf{LiDAR ($905\,\text{nm} / 1550\,\text{nm}$):} Laser pulses scatter off suspended aerosol droplets and rain curtains, producing tens of thousands of phantom return points. These phantom clusters choke point-cloud filtering algorithms, triggering sudden false-positive emergency braking (phantom braking) that causes high-speed rear-end collisions. Furthermore, automotive-grade solid-state and mechanical LiDAR units cost between \$5,000 and \$15,000, prohibiting mass-market integration.
    \item \textbf{Millimeter-Wave Radar ($77\,\text{GHz}$):} While millimeter waves penetrate fog and rain with minimal attenuation, radar sensors possess limited angular resolution and zero vertical elevation discrimination. Consequently, radar cannot detect lane markings, structural asphalt potholes, or pedestrian postures, and frequently misclassifies overhead metal highway signs as stationary obstacles.
    \item \textbf{Monocular CMOS Cameras:} Highly affordable (\$15--\$35) and ubiquitous in commercial fleet dashcams, standard CMOS sensors suffer from severe contrast degradation under fog, sensor dynamic range saturation ($>190$ digital intensity) under oncoming high-beam glare, and excessive thermal shot noise in unlit rural corridors.
    \item \textbf{Long-Wave Infrared (LWIR) Thermal Cameras:} Prohibitively expensive (\$3,000--\$6,000) and highly susceptible to thermal washout during heavy rainfall, where surface evaporative cooling equalizes vehicle and asphalt temperatures.
\end{enumerate}

\subsection{The Edge Compute Dilemma}
Recent breakthroughs in deep convolutional neural networks, vision transformers, and diffusion architectures (e.g., DehazeNet \cite{cai2016dehazenet}, AOD-Net \cite{li2017aodnet}, GridDehazeNet \cite{liu2019griddehazenet}, FFA-Net \cite{qin2020ffa}, DehazeFormer \cite{guo2022dehazeformer}, and DiffDehaze \cite{song2024diffdehaze}) achieve remarkable reconstruction metrics on synthetic benchmark datasets. However, these architectures require high-end discrete GPUs consuming 150--300\,W.

On embedded automotive edge processors (e.g., ARM Cortex-A72, NXP S32G, or NVIDIA Jetson Nano) constrained to strict passive thermal envelopes ($<15\,\text{W}$) and sub-33\,ms deterministic latency deadlines ($\ge 30\,\text{FPS}$), deep restoration models choke severely, exhibiting latencies between 250\,ms and 800\,ms ($<4\,\text{FPS}$). This introduces catastrophic buffer lag, rendering real-time collision prevention impossible.

\subsection{Primary Contributions}
To overcome this fundamental impasse between optical restoration quality, hardware cost, and real-time determinism, this paper introduces \textbf{ClearDrive AI}. The primary contributions are:
\begin{enumerate}
    \item \textbf{Adaptive Foliage-Floor Fast Guided DCP:} An $\mathcal{O}(N)$ linear-time optical dehazing pipeline combining dark channel atmospheric transmission with He's Fast Guided Filter \cite{he2015fastguided}. By introducing a dynamic foliage transmission floor $t_0(x)$ and temporal exponential moving average (EMA) airlight tracking, we resolve the classical ``Tree Bug'' halo artifact and atmospheric sky flickering without iterative optimization.
    \item \textbf{Dual-Domain Optical Dynamic Range Balancing:} Integration of an automotive specular knee compression profile ($>190$ digital intensity rolloff) with luminance-channel Contrast Limited Adaptive Histogram Equalization (CLAHE) \cite{zuiderveld1994clahe} for zero-light Ghost-Vision.
    \item \textbf{Monocular Metric Kinematics \& Collision Warning:} A flat-ground pinhole camera projection coupled with an INT8-quantized YOLOv8n detector \cite{jocher2023yolov8}, deriving metric range $Z$, relative closing velocity $v_{rel}$, dynamic Time-to-Collision ($TTC$), and $+1.5\,\text{s}$ forward ghost-vision cut-in trajectories.
    \item \textbf{Zero-Contention Asynchronous Edge Multiprocessing:} A lock-free shared-memory ring buffer architecture sustaining $32.4\,\text{FPS}$ on an ARM Cortex-A72 processor (Raspberry Pi 4B) with total peak dynamic RAM below $185\,\text{MB}$ and an \$82 BOM.
    \item \textbf{Extensive Field Validation:} Comprehensive empirical validation across standard benchmark datasets (RESIDE \cite{li2018benchmarking}, Foggy Cityscapes \cite{sakaridis2018semantic}) and live mountain highway tests on the 36-hairpin Ooty Ghat road (SH-15), demonstrating a 340\% improvement in forward hazard identification range.
\end{enumerate}

\section{Related Work \& Comparative Survey}

\subsection{Physical Prior and Deep Learning Dehazing}
He et al. \cite{he2010single} formulated the Dark Channel Prior (DCP), observing that in local non-sky patches, the minimum color channel intensity tends toward zero:
\begin{equation}
J^{dark}(x) = \min_{c \in \{r,g,b\}} \left( \min_{y \in \Omega(x)} J^c(y) \right) \approx 0
\label{eq:dcp}
\end{equation}
While revolutionary, He's initial method refined transmission using Levin's Soft Matting Laplacian. Solving the sparse linear system $(L + \lambda U)t = \lambda \tilde{t}$ exhibits $\mathcal{O}(N^2)$ memory and computational complexity, requiring over 18 seconds per 720p frame and crashing embedded systems with Out-Of-Memory (OOM) errors.

He and Sun \cite{he2013guided, he2015fastguided} later introduced the Guided Filter and Fast Guided Filter, achieving $\mathcal{O}(N)$ edge-preserving smoothing. In recent years, deep learning approaches have dominated computer vision literature: DehazeNet \cite{cai2016dehazenet} and AOD-Net \cite{li2017aodnet} reformulated transmission estimation via CNNs; GridDehazeNet \cite{liu2019griddehazenet} and FFA-Net \cite{qin2020ffa} incorporated multi-scale feature attention; while DehazeFormer \cite{guo2022dehazeformer} and DiffDehaze \cite{song2024diffdehaze} applied vision transformers and diffusion priors. Most recently, Chen et al. \cite{chen2024atmospheric} and Zhang et al. \cite{zhang2025realtime} proposed atmospheric scattering-aware networks for autonomous driving. However, as documented by Kumar and Singh \cite{kumar2025ultralow}, deep learning approaches remain computationally intractable for edge microprocessors without dedicated discrete neural accelerators.

\subsection{Low-Light Enhancement \& Specular Glare Control}
Land's Retinex theory \cite{land1977retinex} models an image $S$ as the product of illumination $L$ and reflectance $R$ ($S = L \cdot R$). Multi-scale Retinex with Color Restoration (MSRCR) \cite{jobson1997multiscale} enhances dark regions but causes chromatic drift and noise amplification. Modern approaches like Retinexformer-Edge \cite{wei2024retinexformer} adapt transformer attention for night driving, but single-frame parametric dynamic range compression remains essential for deterministic latency. For glare mitigation, Sharma et al. \cite{sharma2025dynamic} demonstrated that non-linear knee attenuation effectively rolls off high-beam saturation while preserving roadway lane contrast.

\subsection{Real-Time Neural Detection \& Monocular Kinematics}
Real-time object detection was transformed by Redmon et al. \cite{redmon2016yolo} with the You Only Look Once (YOLO) framework. Ultralytics YOLOv8 \cite{jocher2023yolov8}, YOLOv9 \cite{wang2024yolov9}, YOLOv10 \cite{wang2024yolov10}, and YOLO11 \cite{ultralytics2025yolo11} achieved remarkable trade-offs between parameter efficiency and mean Average Precision (mAP). Zhao et al. \cite{zhao2026efficient} and Patel et al. \cite{patel2024monocular} highlighted that projecting bounding-box ground contact points through calibrated pinhole geometry \cite{hartley2003multiple} yields accurate longitudinal distance $Z$ and dynamic Time-to-Collision ($TTC$), fulfilling automotive forward collision warning specifications (AIS-145 \cite{ais145_2020}) without costly multi-camera stereo rigs.

\subsection{In-Cabin Driver Monitoring Systems (DMS)}
Driver fatigue accounts for over 20\% of all commercial vehicular collisions. Soukupov{\'a} and {\v{C}}ech \cite{soukupova2016real} introduced the Eye Aspect Ratio (EAR) for blink classification. Google MediaPipe Face Mesh \cite{kartynnik2019mediapipe} enables real-time 468-point 3D facial topology tracking in under 3.5\,ms on mobile CPUs. Evaluating PERCLOS (Percentage of Eye Closure) \cite{dinges1998perclos} alongside multi-task facial attention \cite{li2024multitask} ensures robust driver drowsiness and distraction alerts.

\begin{table*}[t]
\centering
\caption{Comprehensive Comparative Survey of State-of-the-Art ADAS and Vision Restoration Paradigms}
\label{tab:comparative_survey}
\resizebox{\textwidth}{!}{%
\begin{tabular}{lcccccc}
\toprule
\textbf{System / Architecture} & \textbf{Sensing Modality} & \textbf{Compute Platform} & \textbf{Latency / FPS} & \textbf{Power (W)} & \textbf{Adverse Weather Resilience} & \textbf{Primary Failure Mode} \\
\midrule
Mobileye EyeQ4 (Commercial) & Mono Camera + Radar & EyeQ4 ASIC & 36 FPS & 3--5 W & Moderate (Degrades in dense fog) & Proprietary; closed ecosystem \\
Tesla Vision (HW3/HW4) & Multi-Camera & FSD Dual-NPU & 36 FPS & 72--100 W & Poor (Blinded by high-beam glare/fog) & High power; phantom braking \\
Bosch Front Camera Radar & Radar + Mono Camera & Automotive ECU & 25 FPS & 15--20 W & Good in rain; poor elevation & Cannot detect road potholes \\
DehazeNet (Cai et al. \cite{cai2016dehazenet}) & Monocular Vision & NVIDIA GTX 1080 & 12 FPS & 180 W & Synthetic fog only & Fails in zero-light night \\
AOD-Net (Li et al. \cite{li2017aodnet}) & Monocular Vision & NVIDIA RTX 2080 & 28 FPS & 215 W & Moderate & Low PSNR in dense mountain fog \\
FFA-Net (Qin et al. \cite{qin2020ffa}) & Monocular Vision & NVIDIA V100 GPU & 4.2 FPS & 250 W & High optical quality & Cannot run on edge microprocessors \\
DehazeFormer (Guo \cite{guo2022dehazeformer}) & Monocular Vision & NVIDIA RTX 3090 & 2.1 FPS & 350 W & Exceptional on synthetic sets & Prohibitive latency ($>450\,\text{ms}$) \\
DiffDehaze (Song \cite{song2024diffdehaze}) & Monocular Vision & NVIDIA A100 GPU & 0.8 FPS & 400 W & High realism & Non-deterministic diffusion steps \\
Edge-Guided (Zhang \cite{zhang2025realtime}) & Vision + CAN Bus & Jetson Orin & 24 FPS & 25 W & Good & Expensive hardware (\$600+) \\
\textbf{ClearDrive AI (Proposed)} & \textbf{Monocular Camera} & \textbf{Raspberry Pi 4B} & \textbf{32.4 FPS} & \textbf{7.5 W} & \textbf{Optimal (Fog, Rain, Glare, Night)} & \textbf{Advisory only; no steer actuation} \\
\bottomrule
\end{tabular}%
}
\end{table*}

\section{System Architecture \& Mathematical Formulations}

\begin{figure*}[t]
\centering
\includegraphics[width=0.92\textwidth]{report_figures/Figure_4.1_OmniVision_System_Architecture.png}
\caption{End-to-End Modular Architecture of ClearDrive AI: Video Ingestion, Gyroscopic Pitch Stabilization, OmniVision Dual-Domain Optical Restoration Core, Asynchronous YOLOv8n Kinematics, Metric TTC Telemetry, and Biometric Driver Monitoring System.}
\label{fig:architecture}
\end{figure*}

ClearDrive AI operates as a decoupled, multi-threaded edge perception engine (Fig.~\ref{fig:architecture}). It ingests 720p/1080p video from an onboard CMOS camera or a wireless smartphone dashcam node, stabilizes pitch vibrations, executes optical restoration, runs neural kinematics asynchronously, and projects collision warnings onto a heads-up display.

\subsection{Linear-Time Fast Guided Dark Channel Prior}
Given hazy input image $\mathbf{I}(x)$, the dark channel $I^{dark}(x)$ across local patch $\Omega(x)$ is computed:
\begin{equation}
I^{dark}(x) = \min_{c \in \{r,g,b\}} \left( \min_{y \in \Omega(x)} \frac{I^c(y)}{A^c} \right)
\end{equation}
where atmospheric airlight $\mathbf{A} = (A^r, A^g, A^b)$ is estimated from the top $0.1\%$ brightest pixels in $I^{dark}(x)$. The coarse transmission map $\tilde{t}(x)$ is:
\begin{equation}
\tilde{t}(x) = 1 - \omega \cdot I^{dark}(x)
\label{eq:coarse_trans}
\end{equation}
with haze retention factor $\omega = 0.95$ to preserve depth perception.

To eliminate the $\mathcal{O}(N^2)$ complexity of soft matting, we refine $\tilde{t}(x)$ using the Fast Guided Filter \cite{he2015fastguided}. The guidance image $I$ (grayscale luminance) and input $p = \tilde{t}$ are downsampled by subsampling ratio $s = 4$:
\begin{equation}
I_{sub} = \downarrow_s(I), \quad p_{sub} = \downarrow_s(p)
\end{equation}
The linear coefficients $a_k$ and $b_k$ are evaluated in the subsampled domain using $O(N)$ box filters of radius $r_{sub} = r / s$:
\begin{equation}
\mu_{I} = f_{box}(I_{sub}), \quad \mu_{p} = f_{box}(p_{sub})
\end{equation}
\begin{equation}
\sigma_{I}^2 = f_{box}(I_{sub} \odot I_{sub}) - \mu_{I}^2
\end{equation}
\begin{equation}
\text{cov}_{Ip} = f_{box}(I_{sub} \odot p_{sub}) - \mu_{I} \odot \mu_{p}
\end{equation}
\begin{equation}
a_k = \frac{\text{cov}_{Ip}}{\sigma_{I}^2 + \epsilon}, \quad b_k = \mu_{p} - a_k \odot \mu_{I}
\end{equation}
where $\epsilon = 10^{-3}$ prevents division by zero in flat regions. The coefficient maps are smoothed and upsampled to the native resolution:
\begin{equation}
\bar{a} = \uparrow_s(f_{box}(a_k)), \quad \bar{b} = \uparrow_s(f_{box}(b_k))
\end{equation}
yielding the refined transmission:
\begin{equation}
t_{FGF}(x) = \bar{a}(x) \cdot I(x) + \bar{b}(x)
\label{eq:t_fgf}
\end{equation}
Because box filtering is computed using integral images with 4 array references per pixel regardless of window radius, the time complexity of \eqref{eq:t_fgf} is strictly $\mathcal{O}(N)$, requiring only $1.12\,\text{ms}$ on ARM Cortex-A72 ($16,000\times$ faster than Levin Soft Matting).

\subsection{Eradicating the ``Tree Bug'' and Temporal Sky Flickering}
In practical highway driving tests, classical guided DCP implementations suffer from two critical failure modes:
\begin{enumerate}
    \item \textbf{The Tree Bug (Halo Blooming):} In roadside tree canopies and dark asphalt patches, true scene radiance $J^c(x) \to 0$. From \eqref{eq:coarse_trans}, $\tilde{t}(x)$ drops near zero. Inverting Koschmieder's equation:
    \begin{equation}
    J(x) = \frac{I(x) - A}{t(x)} + A
    \end{equation}
    causes division by a near-zero denominator ($t(x) \to 0$). This forces restored radiance to explode to pure white clipping ($J(x) \to 255$), producing severe white halos around trees and obscuring highway signs.
    \item \textbf{Atmospheric Sky Flickering:} Airlight vector $\mathbf{A}$ computed independently per frame exhibits discrete numerical jitter when oncoming vehicles or overhead cables cross the upper frame boundary, causing severe strobe-like flickering.
\end{enumerate}

\begin{figure}[t]
\centering
\includegraphics[width=\columnwidth]{report_figures/Figure_2.1_DCP_Soft_Matting_vs_Guided_Filter.png}
\caption{Transmission Map Refinement Comparison: (a) Coarse DCP patch map, (b) Levin Soft Matting (18.4\,s), (c) Fast Guided Filter (1.12\,ms), and (d) Final ClearDrive Restoration eliminating foliage halo clipping.}
\label{fig:dcp_comparison}
\end{figure}

ClearDrive AI eliminates these failure modes through two analytical formulations. First, we enforce an adaptive foliage transmission floor $t_0(x)$ inversely proportional to pixel luminance:
\begin{equation}
t_0(x) = \max\left(t_{min}, \, \gamma_{foliage} \cdot \left[1 - Y_{norm}(x)\right]\right)
\label{eq:adaptive_floor}
\end{equation}
where $t_{min} = 0.15$, $\gamma_{foliage} = 0.40$, and $Y_{norm}(x) \in [0, 1]$ is normalized luminance. The bounded transmission map $\hat{t}(x)$ is:
\begin{equation}
\hat{t}(x) = \max\left(t_{FGF}(x), \, t_0(x)\right)
\end{equation}
In dark foliage regions ($Y_{norm} \to 0$), $t_0(x)$ automatically elevates the transmission floor to $0.40$, strictly bounding restored radiance $J(x) \le 255$ and preventing denominator collapse.

Second, temporal sky stability is enforced via an Exponential Moving Average (EMA) filter on airlight:
\begin{equation}
\mathbf{A}_{smooth}^{(t)} = \alpha_{air} \cdot \mathbf{A}^{(t)} + (1 - \alpha_{air}) \cdot \mathbf{A}_{smooth}^{(t-1)}
\label{eq:ema_air}
\end{equation}
with smoothing factor $\alpha_{air} = 0.05$. The complete restored radiance is formulated as:
\begin{equation}
\mathbf{J}(x) = \frac{\mathbf{I}(x) - \mathbf{A}_{smooth}^{(t)}}{\hat{t}(x)} + \mathbf{A}_{smooth}^{(t)}
\label{eq:final_dehaze}
\end{equation}

\subsection{Specular Anti-Glare Knee Compression \& Ghost-Vision HDR}
To counteract blinding high-beam LED headlights, ClearDrive AI implements a continuous piecewise non-linear specular knee compression curve applied to luminance channel $Y(x)$:
\begin{equation}
Y_{out}(x) = 
\begin{cases} 
Y(x), & Y(x) \le Y_{knee} \\
Y_{knee} + \Delta Y \cdot \left(\frac{Y(x) - Y_{knee}}{\Delta Y}\right)^{\gamma_{glare}}, & Y(x) > Y_{knee}
\end{cases}
\label{eq:knee}
\end{equation}
where $Y_{knee} = 190$, $\Delta Y = 255 - Y_{knee} = 65$, and $\gamma_{glare} = 0.35$. At $Y(x) = Y_{knee}$, the curve guarantees $\mathcal{C}^0$ continuity and positive monotonic first derivative:
\begin{equation}
\left. \frac{dY_{out}}{dY} \right|_{Y \to Y_{knee}^+} = \gamma_{glare} \cdot \left(\frac{0}{\Delta Y}\right)^{\gamma_{glare}-1}
\end{equation}
This smoothly attenuates saturated high-beam blooming halos by 68.4\% while leaving darker road surface regions ($Y \le 190$) untouched.

For pitch-black rural corridors ($<0.1\,\text{lux}$), Ghost-Vision HDR applies dual-tile Contrast Limited Adaptive Histogram Equalization (CLAHE) with clip limit $2.0$ across an $8 \times 8$ grid strictly on the luminance channel in YUV color space following black pedestal subtraction ($Y_{floor} = 12$).

\begin{figure}[t]
\centering
\includegraphics[width=\columnwidth]{report_figures/Figure_4.6_Kinematic_Ghost_Vision_Projection.png}
\caption{Kinematic $+1.5\,\text{s}$ Ghost-Vision Predictive Projection: Monocular ground-contact distance $Z$, velocity vector extrapolation, and dynamic Time-to-Collision ($TTC$) calculation.}
\label{fig:ghost_vision}
\end{figure}

\subsection{Monocular Flat-Ground Geometry \& Dynamic TTC}
Obstacle detection is executed by an INT8-quantized YOLOv8n network. For each detected bounding box $\mathbf{B} = [x_1, y_1, x_2, y_2]$, the bottom-center coordinate $(x_{bc}, y_{bottom}) = (\frac{x_1 + x_2}{2}, y_2)$ corresponds to the physical contact point with the roadway.

Assuming flat-ground plane geometry with camera mounting height $H = 1.45\,\text{m}$, vertical focal length $f_y$, and horizon line $y_{horizon}$, longitudinal metric distance $Z$ is computed in closed form:
\begin{equation}
Z = \frac{f_y \cdot H}{\max\left(1.0, \, y_{bottom} - y_{horizon}\right)}
\label{eq:distance}
\end{equation}
Lateral distance $X$ across the vehicle coordinate frame is:
\begin{equation}
X = \frac{(x_{bc} - c_x) \cdot Z}{f_x}
\end{equation}
where $(c_x, c_y)$ is the principal point. To suppress frame-to-frame pixel discretization jitter, raw distances are smoothed via an adaptive filter:
\begin{equation}
Z_{smooth}^{(t)} = \beta_d \cdot Z_{raw}^{(t)} + (1 - \beta_d) \cdot Z_{smooth}^{(t-1)}, \quad \beta_d = 0.70
\end{equation}
The relative closing velocity $v_{rel}$ between host vehicle and target is:
\begin{equation}
v_{rel} = \frac{Z_{smooth}^{(t-1)} - Z_{smooth}^{(t)}}{\Delta t}
\end{equation}
Dynamic Time-to-Collision ($TTC$) is derived:
\begin{equation}
TTC = 
\begin{cases} 
\frac{Z_{smooth}}{v_{rel}}, & v_{rel} > 0.5\,\text{m/s} \\
\infty, & v_{rel} \le 0.5\,\text{m/s}
\end{cases}
\label{eq:ttc}
\end{equation}
In compliance with AIS-145 FCW specifications \cite{ais145_2020}, a two-stage auditory and visual alert is triggered: Critical Warning when $TTC < 1.8\,\text{s}$, and Advisory Warning when $TTC < 2.8\,\text{s}$. Furthermore, dynamic forward trajectories are rendered as $+1.5\,\text{s}$ ``Ghost-Vision'' bounding boxes (Fig.~\ref{fig:ghost_vision}) projected along the target vehicle's kinematic motion vector $\vec{v}_{target}$, warning of sudden cut-in maneuvers before they occur.

\subsection{Biometric Driver Monitoring System (DMS)}
In-cabin driver vigilance is monitored concurrently via Google MediaPipe Face Mesh \cite{kartynnik2019mediapipe}. Eye Aspect Ratio (EAR) is computed from 6 bilateral landmark coordinates:
\begin{equation}
EAR = \frac{\|\mathbf{p}_2 - \mathbf{p}_6\| + \|\mathbf{p}_3 - \mathbf{p}_5\|}{2 \|\mathbf{p}_1 - \mathbf{p}_4\|}
\label{eq:ear}
\end{equation}
A calibrated threshold $EAR_{thresh} = 0.21$ sustained across $N_{consec} = 20$ consecutive frames ($>650\,\text{ms}$) classifies driver microsleep, engaging audible chimes and pulsing cyber-cockpit visual banners.

\begin{algorithm}[t]
\caption{OmniVision Real-Time Restoration Pipeline}
\label{alg:omnivision}
\footnotesize
\begin{algorithmic}[1]
\REQUIRE Input frame $\mathbf{I}$, airlight memory $\mathbf{A}_{smooth}$, parameters $(r, s, \epsilon, \omega, Y_{knee}, \gamma_{glare})$
\ENSURE Restored frame $\mathbf{J}$, updated airlight $\mathbf{A}_{smooth}$
\STATE Convert $\mathbf{I}$ to normalized float32 and extract luminance $Y$
\STATE Compute dark channel $\tilde{I}^{dark}(x) = \min_c \left(\min_{y \in \Omega(x)} \frac{I^c(y)}{A^c}\right)$
\STATE Estimate instantaneous airlight $\mathbf{A}^{(t)}$ from top 0.1\% pixels in $\tilde{I}^{dark}$
\STATE Update temporal airlight: $\mathbf{A}_{smooth} \leftarrow 0.05 \mathbf{A}^{(t)} + 0.95 \mathbf{A}_{smooth}$
\STATE Coarse transmission $\tilde{t}(x) \leftarrow 1 - \omega \tilde{I}^{dark}(x)$
\STATE Downsample $I_{sub} \leftarrow \downarrow_s(Y), \, p_{sub} \leftarrow \downarrow_s(\tilde{t})$
\STATE Compute box filter statistics $\mu_I, \mu_p, \sigma_I^2, \text{cov}_{Ip}$ in subsampled space
\STATE Solve linear regression coefficients $a_k \leftarrow \frac{\text{cov}_{Ip}}{\sigma_I^2 + \epsilon}, \, b_k \leftarrow \mu_p - a_k \mu_I$
\STATE Upsample $\bar{a} \leftarrow \uparrow_s(f_{box}(a_k)), \, \bar{b} \leftarrow \uparrow_s(f_{box}(b_k))$
\STATE Compute refined transmission $t_{FGF}(x) \leftarrow \bar{a}(x) Y(x) + \bar{b}(x)$
\STATE Evaluate foliage floor $t_0(x) \leftarrow \max(0.15, \, 0.40 \cdot [1 - Y(x)/255])$
\STATE Formulate bounded transmission $\hat{t}(x) \leftarrow \max(t_{FGF}(x), \, t_0(x))$
\STATE Restore scene radiance $\mathbf{J}(x) \leftarrow \frac{\mathbf{I}(x) - \mathbf{A}_{smooth}}{\hat{t}(x)} + \mathbf{A}_{smooth}$
\IF{High-Beam Glare Detected}
    \STATE Apply specular knee rolloff on luminance via \eqref{eq:knee}
\ENDIF
\RETURN $\mathbf{J}, \mathbf{A}_{smooth}$
\end{algorithmic}
\end{algorithm}

\begin{figure}[t]
\centering
\includegraphics[width=0.95\columnwidth]{report_figures/Figure_4.2_Async_Worker_Pipeline.png}
\caption{Decoupled Asynchronous Multiprocessing Architecture: Lock-free POSIX shared memory ring buffers isolating sensor ingestion, optical restoration, neural inference, and HTTP/MJPEG streaming.}
\label{fig:multiprocessing}
\end{figure}

\subsection{Asynchronous Multiprocessing Concurrency Model}
As shown in Fig.~\ref{fig:multiprocessing}, standard Python GIL (Global Interpreter Lock) contention causes severe frame dropping when image restoration and deep neural inference run synchronously. ClearDrive AI implements a decoupled multiprocessing architecture using lock-free POSIX shared-memory queues (`latest-frame` dropping policy). 
The perception worker consumes raw frames at $30\,\text{FPS}$, executing the Fast Guided DCP, Anti-Glare, and AR lane projection in $25.8\,\text{ms}$. Neural inference (YOLOv8n and MediaPipe DMS) runs on a secondary worker at $18.2\,\text{ms}$, updating the kinematic telemetry dictionary asynchronously. This decoupling guarantees deterministic display latency ($<28.4\,\text{ms}$) and eliminates buffer bloat.

\begin{algorithm}[t]
\caption{Kinematic Telemetry \& Ghost-Vision Forecasting}
\label{alg:kinematics}
\footnotesize
\begin{algorithmic}[1]
\REQUIRE Detected bounding boxes $\{\mathbf{B}_i\}$, focal length $f_y$, camera height $H$, horizon $y_{horizon}$
\ENSURE Metric telemetry $\{Z_i, v_{rel,i}, TTC_i\}$, Ghost-Vision vectors
\FORALL{detected bounding box $\mathbf{B}_i = [x_{1,i}, y_{1,i}, x_{2,i}, y_{2,i}]$}
    \STATE Ground contact pixel $y_{bottom,i} \leftarrow y_{2,i}$
    \STATE Compute metric distance: $Z_{raw,i} \leftarrow \frac{f_y \cdot H}{\max(1.0, \, y_{bottom,i} - y_{horizon})}$
    \STATE Smooth distance: $Z_{smooth,i}^{(t)} \leftarrow 0.70 Z_{raw,i}^{(t)} + 0.30 Z_{smooth,i}^{(t-1)}$
    \STATE Closing velocity: $v_{rel,i} \leftarrow \frac{Z_{smooth,i}^{(t-1)} - Z_{smooth,i}^{(t)}}{\Delta t}$
    \IF{$v_{rel,i} > 0.5\,\text{m/s}$}
        \STATE Dynamic collision time: $TTC_i \leftarrow \frac{Z_{smooth,i}}{v_{rel,i}}$
    \ELSE
        \STATE $TTC_i \leftarrow \infty$
    \ENDIF
    \STATE Extrapolate $+1.5\,\text{s}$ Ghost-Vision box: $\mathbf{B}_{ghost,i} \leftarrow \mathbf{B}_i + 1.5 \cdot \vec{v}_{target,i}$
    \IF{$TTC_i < 1.8\,\text{s}$}
        \STATE Trigger AIS-145 Critical Audible Chime \& HUD Warning
    \ELSIF{$TTC_i < 2.8\,\text{s}$}
        \STATE Trigger AIS-145 Advisory Visual Alert
    \ENDIF
\ENDFOR
\end{algorithmic}
\end{algorithm}

\section{Experimental Setup and Evaluation}

\subsection{Benchmark Datasets and Testbeds}
The framework was evaluated across both synthetic and real-world datasets:
\begin{itemize}
    \item \textbf{RESIDE Dataset (SOTS Outdoor) \cite{li2018benchmarking}:} 500 outdoor synthetic hazy scenes evaluated for objective optical reconstruction.
    \item \textbf{Foggy Cityscapes \cite{sakaridis2018semantic}:} 550 urban driving scenes with synthetic optical fog densities ($\beta \in \{0.005, 0.01, 0.02\}$).
    \item \textbf{Dark Zurich Benchmark:} Unlit twilight and nighttime automotive scenes for low-light evaluation.
    \item \textbf{ClearDrive Indian Highway Benchmark (CD-IHB):} 12,500 annotated frames captured across real-world adverse road tests: Ooty Mountain Ghat Road (SH-15, 36 hairpin curves, heavy orographic fog), Chennai-Bengaluru Highway (NH-48, nighttime high-beam glare), and torrential monsoon precipitation.
\end{itemize}

\subsection{Hardware Testbeds}
To validate true cross-platform edge deployability, ClearDrive AI was benchmarked on four distinct computing platforms:
\begin{enumerate}
    \item \textbf{Platform A (Embedded ARM SBC):} Raspberry Pi 4 Model B, Quad-core ARM Cortex-A72 @ 1.5\,GHz, 4\,GB LPDDR4, Passive Heatsink.
    \item \textbf{Platform B (Edge AI Compute):} NVIDIA Jetson Nano, Quad-core ARM A57 @ 1.43\,GHz, 128-core Maxwell GPU, 4\,GB LPDDR4.
    \item \textbf{Platform C (Automotive x86 Edge):} Intel Core i5-1135G7 @ 2.40--4.20\,GHz, 8\,GB DDR4.
    \item \textbf{Platform D (Automotive Reference):} NXP BlueBox / S32G Automotive Evaluation Board.
\end{enumerate}

\section{Experimental Results and Discussion}

\begin{figure}[t]
\centering
\includegraphics[width=\columnwidth]{report_figures/Figure_6.1_Throughput_FPS_Comparison.png}
\caption{Throughput (FPS) Comparison across engineering iterations on Intel x86 and Raspberry Pi 4B ARM hardware.}
\label{fig:fps_benchmark}
\end{figure}

\begin{table*}[t]
\centering
\caption{Cross-Iteration Quantitative Benchmark Comparison Across Optical, Throughput, and Memory Metrics}
\label{tab:iteration_comparison}
\resizebox{\textwidth}{!}{%
\begin{tabular}{lcccc}
\toprule
\textbf{Performance Metric} & \textbf{Raw Input} & \textbf{Iteration 1 (Baseline DCP)} & \textbf{Iteration 2 (Fast Guided)} & \textbf{ClearDrive AI (OmniVision)} \\
\midrule
Throughput on x86 CPU (FPS) & 30.0 & 1.8 (Pipeline Choke) & 42.5 & \textbf{38.6} \\
Throughput on ARM Pi 4B (FPS) & 30.0 & 0.0 (OOM Crash) & 22.4 & \textbf{32.4} \\
Transmission Refine Latency (ms) & N/A & 18,420.0 & 1.15 & \textbf{1.12} \\
End-to-End Latency per Frame (ms) & 0.0 & 18,510.0 & 23.5 & \textbf{25.8} \\
Peak Signal-to-Noise Ratio (PSNR, dB) & 12.4 & 21.8 & 20.9 & \textbf{22.6} \\
Structural Similarity Index (SSIM) & 0.58 & 0.84 & 0.81 & \textbf{0.86} \\
Color Difference ($\Delta E_{00}$, CIEDE2000) & 18.4 & 5.2 & 5.8 & \textbf{4.1} \\
Tree Bug Foliage Halo Artifacts & None & Moderate & Severe Blooming & \textbf{Zero Halos (Eliminated)} \\
Atmospheric Sky Flickering & None & High Pulsing & Severe Flickering & \textbf{Zero (EMA Stabilized)} \\
Dynamic System RAM Allocation & 85\,MB & 3,420\,MB & 142\,MB & \textbf{185\,MB} \\
\bottomrule
\end{tabular}%
}
\end{table*}

\begin{table*}[t]
\centering
\caption{State-of-the-Art Quantitative Comparison on RESIDE Outdoor (SOTS) and Foggy Cityscapes Benchmarks}
\label{tab:sota_comparison}
\resizebox{\textwidth}{!}{%
\begin{tabular}{lcccccc}
\toprule
\textbf{Method} & \textbf{Model Type} & \textbf{RESIDE PSNR (dB)} & \textbf{RESIDE SSIM} & \textbf{Cityscapes mAP@0.5} & \textbf{CPU Latency (ms)} & \textbf{Edge FPS (Pi 4B)} \\
\midrule
He et al. (Classical DCP \cite{he2010single}) & Physical Prior & 19.24 & 0.824 & 31.2\% & 18,420.0 & 0.0 (OOM) \\
DehazeNet (Cai et al. \cite{cai2016dehazenet}) & CNN & 21.14 & 0.847 & 35.8\% & 320.0 & 2.8 FPS \\
AOD-Net (Li et al. \cite{li2017aodnet}) & Lightweight CNN & 19.82 & 0.812 & 33.4\% & 145.0 & 6.2 FPS \\
FFA-Net (Qin et al. \cite{qin2020ffa}) & Attention CNN & 23.12 & 0.884 & 38.6\% & 1,420.0 & 0.6 FPS \\
GridDehazeNet (Liu et al. \cite{liu2019griddehazenet}) & Multi-Scale Attention & 23.45 & 0.891 & 39.2\% & 980.0 & 0.9 FPS \\
DehazeFormer (Guo et al. \cite{guo2022dehazeformer}) & Transformer & \textbf{24.18} & \textbf{0.912} & \textbf{41.5\%} & 2,850.0 & 0.2 FPS \\
\textbf{ClearDrive AI (Proposed)} & \textbf{Guided Prior + Knee} & \textbf{22.60} & \textbf{0.860} & \textbf{37.8\%} & \textbf{25.8} & \textbf{32.4 FPS} \\
\bottomrule
\end{tabular}%
}
\end{table*}

\subsection{Quantitative Optical Restoration Analysis}
Table~\ref{tab:iteration_comparison} and Table~\ref{tab:sota_comparison} summarize the progression of ClearDrive AI across three engineering iterations and compare it against state-of-the-art benchmarks. 
Iteration 1 (Classical DCP with Levin Soft Matting) failed catastrophically, requiring 18.4 seconds per frame and crashing on the ARM SBC due to out-of-memory allocation ($>3.4\,\text{GB}$).
Iteration 2 replaced soft matting with He's Fast Guided Filter, reducing refinement latency to $1.15\,\text{ms}$ and boosting throughput to $42.5\,\text{FPS}$. However, as demonstrated in Table~\ref{tab:iteration_comparison}, it suffered from the severe ``Tree Bug'' halo artifact and temporal sky flickering.

ClearDrive AI (Iteration 3) achieved peak optical fidelity with \textbf{22.6\,dB PSNR}, \textbf{0.86 SSIM}, and a low CIEDE2000 color difference of \textbf{4.1}, completely eliminating halos and flickering while locking throughput at \textbf{32.4 FPS} on the ARM Cortex-A72 processor. While deep transformer models (DehazeFormer \cite{guo2022dehazeformer}) achieve slightly higher synthetic PSNR (+1.58\,dB), they require over 2,800\,ms per frame ($0.2\,\text{FPS}$ on edge CPU), representing a $110\times$ latency penalty that violates automotive safety requirements.

\begin{figure}[t]
\centering
\includegraphics[width=\columnwidth]{report_figures/Figure_6.2_PSNR_SSIM_Quality_Restoration.png}
\caption{Reconstruction Fidelity (PSNR and SSIM) across varying fog density levels on the RESIDE and Foggy Cityscapes benchmarks.}
\label{fig:psnr_ssim}
\end{figure}

\subsection{Cross-Platform Compute Benchmarks}
Table~\ref{tab:hardware_benchmarks} details hardware utilization across platforms. ClearDrive AI sustains deterministic real-time operation across all architectures:
\begin{itemize}
    \item On the \$82 Raspberry Pi 4B, the system maintains \textbf{32.4 FPS} at 74\% CPU load and $178\,\text{MB}$ RAM, operating safely within thermal limits ($62^\circ\text{C}$).
    \item On the NVIDIA Jetson Nano, throughput reaches \textbf{34.8 FPS} utilizing GPU tensor acceleration.
    \item On modern x86 hardware, throughput exceeds \textbf{38.6 FPS} at only 38\% CPU utilization.
\end{itemize}

\begin{table}[t]
\centering
\caption{Cross-Platform Edge Hardware Benchmarks}
\label{tab:hardware_benchmarks}
\resizebox{\columnwidth}{!}{%
\begin{tabular}{lcccc}
\toprule
\textbf{Hardware Testbed} & \textbf{Pipeline FPS} & \textbf{CPU Load} & \textbf{RAM} & \textbf{Temp} \\
\midrule
Intel Core i5-1135G7 & 38.6 & 38\% & 185\,MB & 52$^\circ$C \\
NVIDIA Jetson Nano & 34.8 & 54\% / 62\% & 215\,MB & 58$^\circ$C \\
\textbf{Raspberry Pi 4B (ARM)} & \textbf{32.4} & \textbf{74\%} & \textbf{178\,MB} & \textbf{62$^\circ$C} \\
NXP BlueBox Reference & 31.5 & 68\% & 192\,MB & 55$^\circ$C \\
\bottomrule
\end{tabular}%
}
\end{table}

\subsection{Real-World Highway Field Trials}
Road tests on the Nilgiri Mountain Ghat Road (SH-15, Ooty, India) demonstrated exceptional hazard mitigation:
\begin{enumerate}
    \item \textbf{Mountain Fog Hazard Identification Range:} Under dense advection fog reducing human visibility to $<10\,\text{m}$, baseline YOLOv8n failed to detect an oncoming passenger bus until it reached $8.2\,\text{m}$ (providing an impossible $0.98\,\text{s}$ driver reaction window). With ClearDrive AI optical restoration active, detection locked reliably onto the bus at \textbf{38.5 meters} ($4.6\,\text{s}$ reaction window)---yielding a \textbf{340\% improvement} in hazard warning distance.
    \item \textbf{Monsoon Torrential Rain:} Atmospheric spray was suppressed while virtual AR lane guides prevented vehicle runoff during lane-marker washouts.
    \item \textbf{Night Specular Glare Suppression:} High-beam LED saturation was compressed by 68.4\%, preserving drivable lane edges without dark-region clipping.
\end{enumerate}

\begin{figure}[t]
\centering
\includegraphics[width=\columnwidth]{report_figures/Figure_6.3_Distance_Calibration_Curve.png}
\caption{Monocular Pinhole Ground-Plane Metric Distance Calibration Curve: Predicted vs. Ground-Truth LiDAR Distance up to $60\,\text{m}$.}
\label{fig:distance_calibration}
\end{figure}

\subsection{Monocular Distance Accuracy}
Fig.~\ref{fig:distance_calibration} illustrates the distance calibration curve verified against physical laser rangefinders. From $2\,\text{m}$ to $25\,\text{m}$, metric error remained under $\pm 3.8\%$. At extended distances ($25\,\text{m}$ to $50\,\text{m}$), error remained bounded within $\pm 7.2\%$, fully satisfying AIS-145 FCW alert safety criteria.

\subsection{Ablation Study}
To evaluate the individual contributions of each algorithmic component, we performed an extensive ablation study (Table~\ref{tab:ablation}).
Removing the Fast Guided Filter (falling back to coarse patch boxes) caused severe step-edge artifacts (PSNR dropped from $22.6\,\text{dB}$ to $16.8\,\text{dB}$). Removing the temporal EMA airlight induced disorienting strobe flickering. Omission of the adaptive foliage floor $t_0(x)$ caused the Tree Bug to reappear, creating blinding white halos around overhead signs. Finally, executing without asynchronous multiprocessing caused pipeline throughput to collapse to $8.4\,\text{FPS}$, demonstrating the indispensability of the decoupled shared-memory architecture.

\begin{table}[t]
\centering
\caption{Ablation Study of Algorithmic Components}
\label{tab:ablation}
\resizebox{\columnwidth}{!}{%
\begin{tabular}{lccc}
\toprule
\textbf{Configuration} & \textbf{FPS} & \textbf{PSNR / SSIM} & \textbf{Observed Impact / Failure Mode} \\
\midrule
\textbf{Full ClearDrive AI} & \textbf{32.4} & \textbf{22.6 dB / 0.86} & \textbf{Optimal; zero halos; steady telemetry} \\
w/o Fast Guided Filter & 34.1 & 16.8 dB / 0.71 & Step-edge artifacts; silhouette bleed \\
w/o Temporal EMA Airlight & 32.6 & 20.1 dB / 0.79 & Violent strobe-like sky flickering \\
w/o Foliage Floor $t_0(x)$ & 32.5 & 19.2 dB / 0.77 & Tree Bug returns: white canopy halos \\
w/o Gyro Pitch Dampener & 32.8 & 21.5 dB / 0.83 & Vertical road bounce (TTC err $\pm 38\%$) \\
w/o Async Multiprocessing & 8.4 & 22.4 dB / 0.85 & Throughput collapse; buffer lag ($>350\,\text{ms}$) \\
\bottomrule
\end{tabular}%
}
\end{table}

\section{Discussion, Functional Safety \& Economic Impact}

\subsection{Fail-Safe Architectural Compliance}
ClearDrive AI is designed as an advisory, human-in-the-loop ADAS operating strictly within ISO 26262 ASIL-B \cite{iso26262_2024} and AIS-140/145 guidelines. To prevent sensor lockup hazards:
\begin{enumerate}
    \item \textbf{Hardware Bypass Toggle:} A single-cycle hardware bypass toggle allows immediate $1:1$ routing of raw dashcam frames ($<15\,\text{ms}$ fail-safe handover).
    \item \textbf{Offline-First Edge Operation:} Zero dependency on cloud APIs ensures 100\% operational availability across remote mountain tunnels and zero-cellular dead zones, preventing telemetry drops.
\end{enumerate}

\subsection{Economic Feasibility and Democratic Impact}
Commercial automotive ADAS packages (e.g., Mobileye EyeQ4, Bosch front-camera radar suites) retail between \$1,500 and \$4,500, restricting life-saving collision prevention to luxury vehicle tiers. ClearDrive AI achieves comparable perception robustness on an $\$82$ (₹6,800) hardware bill of materials (Raspberry Pi 4B: \$58, Sony IMX323 CMOS sensor: \$18, and automotive 12V buck regulator: \$6), enabling mass-market adoption in developing economies.

\section{Conclusion and Future Work}
This paper presented \textbf{ClearDrive AI}, a real-time, deterministic, and edge-resilient multi-hazard vision restoration and kinematic telemetry ADAS framework. By coupling an adaptive foliage-floor Fast Guided Dark Channel Prior with specular knee dynamic range compression and asynchronous multiprocessing, ClearDrive AI permanently resolves the classical halo and sky-flickering failure modes of physical dehazing, sustaining \textbf{32.4 FPS} on an \$82 ARM Cortex-A72 edge SBC. Field validation across standard benchmarks and 36-hairpin mountain ghat road tests demonstrates an expansion in forward hazard identification range by 340\% in dense fog, with bounded metric distance error under 7.2\% and dynamic TTC alert generation.

Future research will focus on extending the monocular geometric pipeline to stereo-vision pseudo-LIDAR fusion and deploying 4-bit INT4 quantization via TensorRT on ultra-low-power NPU micro-controllers.

\section*{Acknowledgment}
The authors express sincere gratitude to the Department of CSE-AIML, the faculty guides, and the research leadership at Chennai Institute of Technology, Chennai, India, for providing laboratory facilities, technical mentorship, and computational resources.

\bibliographystyle{IEEEtran}
\bibliography{references}

\end{document}
